\documentclass[11pt,a4paper]{article}
\usepackage{webclases-lesson}
\lessonSetup{T5}{v0.2}{T5 - Energia y trabajo}
\begin{document}
\shorthandoff{<>}
\color{Ink}
\coverHeader{T5 --- Energ\'ia y trabajo}{Comparar estados, reconocer transferencias y construir balances.}
El trabajo es una transferencia de energ\'ia mediante fuerzas.
La energ\'ia es una magnitud escalar; un balance permite comparar estados sin
calcular siempre todos los instantes del movimiento.
\begin{lessonindex}
\indexentry{sec:b1}{1. Trabajo, rapidez y potencia}
\indexentry{sec:b2}{2. Energ\'ia almacenada en interacciones}
\indexentry{sec:b3}{3. Balances y etapas}
\end{lessonindex}
Usaremos $g=10\mpss$ y el modelo de part\'icula de masa constante.
\precisionnote

\block{1}{Trabajo, rapidez y potencia}{sec:b1}
Para fuerza constante, el producto escalar de T0 da
\[
W=\vec F\cdot\Delta\vec r=F|\Delta\vec r|\cos\theta.
\]
El \'angulo es el formado por la fuerza y el desplazamiento, no necesariamente
el \'angulo de la rampa. El trabajo es positivo si la componente de fuerza
acompa\~na al desplazamiento, negativo si se opone y cero si es perpendicular.
En un tramo recto sin inversi\'on, $|\Delta\vec r|=d$.
En una trayectoria general:
\[
W=\int\vec F\cdot d\vec r,\qquad
W=\int_{x_i}^{x_f}F_x(x)\,dx\quad\text{en una dimensi\'on}.
\]
Una fuerza grande puede hacer trabajo cero si no hay desplazamiento o si es
siempre perpendicular a \'el. Por ejemplo, la resultante radial del MCU cambia
la direcci\'on, pero no la rapidez.

\heading{El trabajo total cambia la energ\'ia cin\'etica}
La energ\'ia cin\'etica es $K=mv^2/2$. En un tramo rectil\'ineo,
Newton y la regla de la cadena dan
$F_x=ma_x=mv\,dv/dx$; al integrar:
\[
W_{\mathrm{total}}=\int mv\,dv
=\frac12mv_f^2-\frac12mv_i^2=\Delta K.
\]
El teorema cuenta el trabajo de \emph{todas} las fuerzas.
Un trabajo individual negativo no obliga a perder rapidez si otro trabajo positivo
lo supera. $K$ no tiene signo de direcci\'on: si buscamos la velocidad, su sentido
se decide aparte.

\begin{examplebox}{empujar hacia abajo o tirar hacia arriba}
\origin{Adaptado del original T5, pp.\ 4--5.}
\step{Planteamiento} Bloque de $1,00\kg$ que ya desliza hacia la derecha a
$3,00\mps$, $F=8,00\N$, $d=6,00\units{m}$ y $\mu_k=0,30$.
Compara fuerza a $30^\circ$ por debajo y por encima de la horizontal:
trabajo de cada fuerza y rapidez final.
\solutionblock
\step{Modelo} No hay aceleraci\'on vertical. Empujar hacia abajo aumenta $N$
y el rozamiento; tirar hacia arriba los reduce. El trabajo aplicado ser\'a
igual porque $F_x$ es la misma.
\begin{center}
\begin{tikzpicture}[font=\small]
\foreach \shift/\angle/\lab in {0/-30/empujar,6/30/tirar} {
 \begin{scope}[xshift=\shift cm]
 \draw[fill=BlueSoft,draw=BrandBlue] (-.35,-.3) rectangle (.35,.3);
 \draw[force] (0,0)--++(\angle:1.8) node[above right] {$F$};
 \draw[axis] (-.3,-1)--(2,-1) node[right] {$\Delta x$};
 \draw[guide] (0,0)--(1.7,0);
 \node at (1,.35) {$30^\circ$};
 \node at (0,1.2) {\lab};
 \end{scope}
}
\end{tikzpicture}
\end{center}
Esquema de fuerza aplicada y desplazamiento; las restantes fuerzas se cuentan abajo.
\step{C\'alculo} Al empujar, $N=mg+F\sin30^\circ=14,0\N$ y $f_k=4,20\N$.
\[
W_F=Fd\cos30^\circ=41,6\J,\quad
W_f=-f_kd=-25,2\J,\quad W_N=W_P=0.
\]
Con $K_i=4,50\J$, $K_f=20,869\ldots\J$ y $v_f=6,46\mps$.
Al tirar, $N=mg-F\sin30^\circ=6,00\N>0$,
$f_k=1,80\N$ y $W_f=-10,8\J$.
Resultan $K_f=35,269\ldots\J$ y $v_f=8,40\mps$.
\step{Interpretaci\'on} La componente vertical no hace trabajo directamente,
pero modifica el rozamiento. En ambos casos se conserva el contacto y el
movimiento sigue a la derecha: las hip\'otesis usadas son coherentes.
\end{examplebox}

\begin{examplebox}{fuerza variable y m\'aximo de rapidez}
\origin{Elaboraci\'on propia.}
\step{Planteamiento} Bloque de $2,00\kg$ en reposo en $x=0$, sobre mesa lisa.
La fuerza neta es $F_x=6\units{N}-(2\units{N/m})x$ hasta $x=4,00\units{m}$.
Halla d\'onde alcanza la mayor rapidez y cu\'al es al final.
\solutionblock
\step{Modelo} No hay aceleraci\'on constante, porque cambia la fuerza.
Integramos $F(x)$, o calculamos sus \'areas con signo, para hallar $\Delta K$.
\begin{center}
\begin{tikzpicture}[x=1.25cm,y=.42cm,font=\small]
\fill[GreenLine!20] (0,0)--(0,6)--(3,0)--cycle;
\fill[AmberLine!20] (3,0)--(4,-2)--(4,0)--cycle;
\draw[axis] (0,0)--(4.8,0) node[right] {$x$ (m)};
\draw[axis] (0,-2.8)--(0,7.4) node[above] {$F_x$ (N)};
\draw[BrandBlue,thick,domain=0:4] plot(\x,{6-2*\x});
\foreach \x in {1,2,3,4} \draw (\x,0)--(\x,-.15) node[below] {$\x$};
\foreach \y in {-2,2,4,6} \draw (0,\y)--(-.05,\y) node[left] {$\y$};
\node at (1,2) {$+9\J$}; \node at (3.55,-1.25) {$-1\J$};
\end{tikzpicture}
\end{center}
\step{C\'alculo}
\[
W(0\to x)=(6\units{N})x-(1\units{N/m})x^2.
\]
En $x=3,00\units{m}$ cambia el signo de $F$:
$K=9,00\J$ y $v_{\max}=\sqrt{2K/m}=3,00\mps$.
De $3$ a $4\units{m}$ el trabajo es $-1,00\J$, luego
$K_f=8,00\J$ y $v_f=2,83\mps$.
\step{Interpretaci\'on} En $x=3$ la fuerza es cero, pero la velocidad no.
El bloque sigue avanzando mientras pierde parte de la energ\'ia ganada.
\variation{Si la misma ley se prolongase, $K=6x-x^2=0$ de nuevo en
$x=6,00\units{m}$. Este balance localiza la siguiente parada, pero no calcula su tiempo.}
\end{examplebox}

\heading{Potencia: cu\'an r\'apido se transfiere energ\'ia}
La potencia media es $P_{\mathrm{med}}=W/\Delta t$ y para una fuerza,
$P=\vec F\cdot\vec v$ instant\'aneamente.
El julio (J) mide energ\'ia y el vatio (W) energ\'ia por segundo.
$1\units{kWh}=3,6\units{MJ}$ es una energ\'ia, no una potencia.
Con rendimiento $\eta$, $P_{\mathrm{util}}=\eta P_{\mathrm{entrada}}$.

\begin{examplebox}{potencia de un elevador}
\origin{Elaboraci\'on propia; sustituye la comparaci\'on de bombillas del original.}
\step{Planteamiento} Se suben $50,0\kg$ a rapidez constante $0,400\mps$
hasta $4,00\units{m}$, con $\eta=0,80$. Halla potencia \'util y energ\'ia de entrada.
\solutionblock
\step{Modelo} $a=0$ exige fuerza ascendente $mg$. Su potencia es $mgv$.
\step{C\'alculo} $P_{\mathrm{util}}=200\units{W}$,
$P_{\mathrm{entrada}}=200/0,80=250\units{W}$.
$t=h/v=10,0\units{s}$ y $E_{\mathrm{entrada}}=250(10)=2,50\units{kJ}$.
\step{Interpretaci\'on} Solo $mgh=2,00\units{kJ}$ aumentan la energ\'ia potencial.
Si se eleva a doble rapidez con el mismo rendimiento, se duplica la potencia,
se divide el tiempo por dos y se mantiene la energ\'ia de entrada.
\end{examplebox}

\block{2}{Energ\'ia almacenada en interacciones}{sec:b2}
Una fuerza conservativa tiene trabajo independiente del camino entre dos estados.
Podemos asociarle una energ\'ia potencial $U$ tal que $W_{\mathrm{cons}}=-\Delta U$.
La potencial corresponde a una interacci\'on del sistema, como objeto--Tierra;
no es una nueva fuerza ni un trabajo acumulado dentro de un objeto aislado.

\heading{Gravedad y elecci\'on del cero}
En gravedad uniforme, $U_g=mgh$. El cero de altura se elige por conveniencia.
$U_g$ puede ser negativa si $h<0$; eso no impide el movimiento.
Lo que interviene en el balance es $\Delta U_g$.
Al bajar, el peso realiza trabajo positivo y $U_g$ disminuye.

\begin{examplebox}{subir y volver con la misma energ\'ia}
\origin{Adaptado del original T5, p.\ 9.}
\step{Planteamiento} Una pelota de $0,500\kg$ sale de $h=0$ a $+12,0\mps$.
Sin aire, halla rapidez a $5,00\units{m}$ y altura m\'axima.
\solutionblock
\step{Modelo} En el sistema pelota--Tierra, solo cambia la distribuci\'on
entre $K$ y $U_g$, por lo que $K+U_g$ es constante.
\step{C\'alculo} $K_i=36,0\J$.
A $h=5,00\units{m}$, $U_g=25,0\J$, $K=11,0\J$ y
$|\vec v|=\sqrt{2(11)/0,500}=6,63\mps$.
En el m\'aximo $K=0$: $h_{\max}=36/(0,500\cdot10)=7,20\units{m}$.
\step{Interpretaci\'on} A $5\units{m}$ puede tener $v_y=+6,63\mps$ al subir
o $-6,63\mps$ al bajar. El balance determina rapidez, no sentido.
Al volver a $h=0$, $v_y=-12,0\mps$.
\variation{Elegir el cero dos metros m\'as abajo suma $10,0\J$ a $U_g$ y a la
energ\'ia mec\'anica de todos los estados; no cambia ninguna velocidad.}
\end{examplebox}

\heading{Energ\'ia de un muelle}
Para deformar lentamente un muelle ideal desde $0$ hasta $x$ se aplica una fuerza
de m\'odulo creciente $k|x|$. El \'area triangular bajo esa fuerza da
\[
U_{\mathrm{el}}=\frac12kx^2.
\]
$x$ se mide desde la longitud natural, tambi\'en al comprimir.
En un oscilador horizontal ideal unido al muelle, con amplitud $A$:
$E=kA^2/2$ y $K=k(A^2-x^2)/2$ para $|x|\le A$.
$K$ se anula en ambos extremos, $+A$ y $-A$; en $x=0$ es m\'axima.
La velocidad puede tener cualquiera de los dos signos seg\'un la etapa.
No estamos introduciendo un modelo completo de oscilaci\'on temporal.

\block{3}{Balances y etapas}{sec:b3}
El balance que utilizaremos es
\[
\Delta(K+U)=W_{\mathrm{otras}}.
\]
Las fuerzas incluidas mediante potencial no vuelven a contarse como trabajos en
el segundo miembro. Si solo act\'uan esas fuerzas, $K+U$ se conserva.
En un bloque deslizante sobre una superficie fija con $f_k$ constante,
$W_f=-f_kd$: parte de la energ\'ia mec\'anica pasa a energ\'ia interna.
La energ\'ia total no desaparece.

\begin{examplebox}{ca\'ida con resistencia idealizada}
\origin{Adaptado del original T5, pp.\ 10--11.}
\step{Planteamiento} Cuerpo de $1,00\kg$ desde reposo a $10,0\units{m}$.
Compara ca\'ida libre y resistencia constante de $2,00\N$ hacia arriba.
Halla rapidez al llegar al suelo y balance a $h=4,00\units{m}$.
\solutionblock
\step{Modelo} Para objeto--Tierra, $E_i=100\J$ y
$W_r=-2(10-h)$. La resistencia constante es una simplificaci\'on,
no una ley general del aire.
\step{C\'alculo} A $h=4$, $U_g=40,0\J$.
Sin resistencia $K=60,0\J$; con ella $K=100-40-12=48,0\J$.
Al suelo, $K=100\J$ sin aire y $80,0\J$ con resistencia:
rapideces $14,1$ y $12,6\mps$, respectivamente.
\begin{center}
\begin{tikzpicture}[x=1cm,y=.032cm,font=\small]
\draw[axis] (0,0)--(0,120) node[above] {energ\'ia (J)};
\foreach \e in {0,40,80,100} \draw[guide] (0,\e)--(8,\e) node[right] {$\e$};
\draw[fill=BlueSoft,draw=BrandBlue] (1,0) rectangle (2,100);
\draw[fill=GreenSoft,draw=GreenLine] (3.5,0) rectangle (4.5,80);
\draw[fill=AmberSoft,draw=AmberLine] (6,0) rectangle (7,20);
\node[below,align=center] at (1.5,0) {inicio\\$U_g=100$};
\node[below,align=center] at (4,0) {suelo\\$K=80$};
\node[below,align=center] at (6.5,0) {transferida\\$E_{\mathrm{dis}}=20$};
\end{tikzpicture}
\end{center}
\step{Interpretaci\'on} Las barras usan una escala com\'un:
$100=80+20\J$. La energ\'ia disipada indica transferencia a energ\'ia interna
y entorno, no una tercera energ\'ia mec\'anica del cuerpo.
No hay una barra de trabajo como energ\'ia almacenada final.
\end{examplebox}

\begin{examplebox}{del muelle a la rampa}
\origin{Adaptado del original T5, p.\ 14.}
\step{Planteamiento} Bloque de $0,500\kg$ inicialmente en reposo contra un muelle
horizontal de $100\units{N/m}$, comprimido $0,200\units{m}$.
No est\'a unido a \'el; sale a una rampa lisa mediante una transici\'on suave.
Calcula rapidez de salida y altura m\'axima.
\solutionblock
\step{Modelo} Sistema bloque--muelle--Tierra, sin rozamiento.
Primero $U_{\mathrm{el}}$ pasa a $K$; despu\'es $K$ pasa a $U_g$.
La salida del muelle ocurre en su longitud natural.
\step{C\'alculo}
\[
E_i=\tfrac12(100)(0,200)^2=2,00\J,\quad
v_{\mathrm{salida}}=\sqrt{\frac{2E_i}{m}}=2,83\mps,\quad
h_{\max}=\frac{E_i}{mg}=0,400\units{m}.
\]
\step{Interpretaci\'on} Rampas lisas de $30^\circ$ y $60^\circ$ permiten la misma
altura. Si idealizamos el inicio del tramo recto al nivel de referencia,
$s=h/\sin\theta$ da $0,800$ y $0,462\units{m}$: cambia el recorrido, no el balance.
\variation{Si antes de la rampa hay $0,500\units{m}$ rugosos con $\mu_k=0,20$,
$W_f=-0,500\J$. Queda $1,50\J>0$, as\'i que alcanza la rampa y sube
$0,300\units{m}$. Antes de usar una altura final hay que comprobar que puede llegar.}
\end{examplebox}

\heading{El sistema cambia la escritura, no el resultado}
Un bloque sube una rampa lisa con $K_i=4,00\J$; una cuerda hace $+18,0\J$
y la potencial gravitatoria aumenta $12,0\J$.
Para bloque--Tierra: $\Delta(K+U_g)=18$ y $K_f=10,0\J$.
Para el bloque solo: $\Delta K=W_{\mathrm{cuerda}}+W_P=18-12=6,00\J$.
Ambos balances coinciden. Contar $W_P$ adem\'as de $\Delta U_g$ en el primero
descontar\'ia dos veces la misma interacci\'on.

\begin{exercisebox}{subir sin ganar rapidez}
Se eleva a velocidad constante un objeto de $2,00\kg$ una altura de
$3,00\units{m}$. Calcula trabajo aplicado, trabajo del peso y $\Delta K$.
Explica por qu\'e aumenta $U_g$ aunque no aumente $K$.
\solutionblock
$W_{\mathrm{aplicada}}=+60,0\J$, $W_P=-60,0\J$ y
$\Delta K=W_{\mathrm{total}}=0$.
Para objeto--Tierra, el trabajo externo aumenta $U_g$ en $60,0\J$.
La potencia depender\'ia adem\'as del tiempo empleado.
\end{exercisebox}

\begin{warnbox}
$W_f=-f_kd$ corresponde aqu\'i a deslizamiento sobre una superficie fija.
No se generaliza sin estudiar el contacto a cintas m\'oviles o rodadura.
En procesos por etapas, cambia de ecuaci\'on donde cambia el modelo y conserva
el estado final como inicio de la etapa siguiente.
\end{warnbox}

\heading{Fuentes y procedencia}
Empuje inclinado, tiro vertical, ca\'ida resistente y muelle--rampa adaptan
el original T5 y su background. Fuerza variable, elevador, comparaci\'on de
sistemas y actividades son de elaboraci\'on propia. Diagramas y cuentas propios.
\source{https://fisquiweb.es/Apuntes/apun1B.htm}{FisQuiWeb / IES La Magdalena: inspiraci\'on principal, 1.\textordmasculine\ Bachillerato.}
Contraste externo: alcance del teorema del trabajo total, no del trabajo de una fuerza aislada.
\source{https://openstax.org/books/university-physics-volume-1/pages/7-3-work-energy-theorem}{OpenStax, University Physics 1, 7.3: Work-Energy Theorem.}
\end{document}
